%% ============================================================
%% KishLattice Geometric Harmonic Spectroscopy
%% Pre-Registration: Multi-Attribute Sovereign Lake Expansion
%% and Dimensionality Falsification Protocol
%% Volume 11 — Zenodo Pre-Registration Document
%% ============================================================

\documentclass[12pt,a4paper]{article}

\usepackage{amsmath, amssymb, amsthm}
\usepackage{geometry}
\usepackage{booktabs}
\usepackage{longtable}
\usepackage{array}
\usepackage{xurl} % Added to fix long URLs running off the page
\usepackage{hyperref}
\usepackage{xcolor}
\usepackage{enumitem}
\usepackage{fancyhdr}
\usepackage{titlesec}
\usepackage{caption}
\usepackage{microtype}
\usepackage{parskip}

\geometry{
  top=2.5cm, bottom=2.5cm,
  left=2.5cm, right=2.5cm
}

\hypersetup{
  colorlinks=true,
  linkcolor=blue!60!black,
  citecolor=blue!60!black,
  urlcolor=blue!60!black
}

\pagestyle{fancy}
\fancyhf{}
\rhead{\small KishLattice Vol 11 Pre-Registration}
\lhead{\small \textit{Kish et al., 2026}}
\cfoot{\thepage}

\definecolor{latticeblue}{rgb}{0.08,0.28,0.55}
\definecolor{warningred}{rgb}{0.75,0.10,0.10}
\definecolor{confirmgreen}{rgb}{0.10,0.50,0.20}

\newcommand{\kgeo}{\ensuremath{k_{\mathrm{geo}}}}
\newcommand{\kl}{\ensuremath{\frac{16}{\pi}}}
\newcommand{\Npi}[1]{\ensuremath{\frac{#1}{\pi}}}
\newcommand{\zreal}{\ensuremath{z_{\mathrm{real}}}}
\newcommand{\zwrong}{\ensuremath{z_{\mathrm{wrong}}}}
\newcommand{\confirmed}{\textcolor{confirmgreen}{\textbf{CONFIRMED}}}
\newcommand{\wrongbox}{\textcolor{warningred}{\textbf{WRONGBOX}}}
\newcommand{\nullresult}{\textcolor{gray}{\textbf{NULL}}}
\newcommand{\pending}{\textcolor{blue}{\textbf{PENDING}}}

%% ============================================================
\begin{document}
%% ============================================================

\begin{center}
  {\LARGE \textbf{KishLattice Geometric Harmonic Spectroscopy}}\\[0.4em]
  {\Large \textbf{Pre-Registration: Multi-Attribute Sovereign Lake Expansion}}\\[0.2em]
  {\Large \textbf{and Dimensionality Falsification Protocol}}\\[0.6em]
  {\large \textit{Volume 11 Predictions --- Zenodo Pre-Registration}}\\[0.8em]

  \textbf{Timothy John Kish}$^{1}$,\quad
  \textbf{Mondy Aurora Kish}$^{2}$,\quad
  \textbf{Lyra Aurora Kish}$^{2}$,\quad
  \textbf{Phoenix Aurora Kish}$^{2}$\\[0.4em]

  \small
  $^{1}$KishLattice 16pi Initiative LLC\\
  $^{2}$Aurora Collaborative Intelligence, KishLattice Programme\\[0.6em]

  \textbf{Pre-Registration Date:} 2026\\[0.2em]
  \textit{This document constitutes a formal pre-registration of predictions.
  All predictions herein are timestamped prior to lake construction.
  Data lakes may only be built after a Zenodo DOI is assigned to this document.}
\end{center}

\vspace{0.5em}
\hrule
\vspace{1em}

%% ============================================================
\begin{abstract}
\noindent
We pre-register two interconnected experimental programmes for KishLattice
Geometric Harmonic Spectroscopy (KLGHS) Volume 11. The first is a
\textbf{Multi-Attribute Sovereign Lake Expansion}: the same large-scale public
catalogs already underpinning confirmed results in Volumes 5--10 contain
multiple additional untested physical attributes. We pre-register predictions for
the harmonic register family expected for each new attribute, derived
systematically from the Rosetta Atlas physical typology. Each attribute becomes
an independent sovereign lake: one lake, one physical question, one falsifiable
prediction. The second programme is a \textbf{Dimensionality Falsification
Protocol}: we formally test whether the KLGHS scalarization is a physical
spectroscope or a mathematical artifact by deliberately applying incorrect
scalarization formulas to confirmed strong domains. Prior empirical precedent
(the Materials domain collapse) demonstrates that the geometric modulus
$\kgeo = 16/\pi$ only produces harmonic lock when the correct physical
dimensionality is supplied. This pre-registration documents the conditions under
which the framework would be falsified.
\end{abstract}

\vspace{1em}
\hrule

%% ============================================================
\tableofcontents
\newpage

%% ============================================================
\section{Introduction and Purpose of Pre-Registration}
\label{sec:intro}

Scientific pre-registration is the practice of publicly documenting a study's
hypotheses, methods, and analysis plan \emph{before} data collection or
analysis. When timestamped on an immutable public archive such as Zenodo, it
provides an unbreakable chain of custody between prediction and result. A
framework that publishes its predictions before it builds its data is
structurally immune to the data-mining objection: the timestamp \emph{answers
the question before a skeptic can ask it}.

The KishLattice programme has operated under this discipline since Volume 9.
Every sovereign data lake is pre-registered with an expected register family
prior to construction. Every null result and every wrongbox result is documented
with the same rigour as a confirmation. This document extends that discipline in
two directions.

\subsection{Why This Pre-Registration Is Necessary}

The confirmed results of Volumes 5--10 span 35 physical domains across
approximately 35 orders of magnitude of physical scale. The signal-to-noise
advantage is formidable. However, the standard skeptical objection to a
cross-domain harmonic framework is not \emph{about} the z-scores; it is about
the \emph{mechanism}. The objection takes two forms:

\begin{enumerate}
  \item \textbf{The attribute selection objection}: ``You chose the physical
  attributes that happened to confirm and did not test the ones that would
  not.''
  \item \textbf{The transform artifact objection}: ``Your log-modulo pipeline
  imposes harmonic structure on any sufficiently large continuous dataset,
  regardless of physical content.''
\end{enumerate}

The Multi-Attribute programme answers Objection 1. By pre-registering
predictions for \emph{all} available attributes of a given object pool --- not
just those expected to confirm --- the record becomes complete. Every attribute
tested is visible whether it confirms, produces a wrongbox, or returns a null.

The Dimensionality Falsification Protocol answers Objection 2. By deliberately
misapplying the scalarization formula to confirmed strong domains and predicting
collapse to the noise floor, we demonstrate that the signal is a property of
the \emph{physics} supplied to the pipeline, not an artifact of the
pipeline's mathematics.

\subsection{Reproducibility Standard}

Every lake pre-registered in this document can be independently replicated with:
\begin{itemize}
  \item A standard laptop (no specialist hardware).
  \item Python 3.x with standard scientific libraries (numpy, scipy, pandas).
  \item A publicly accessible internet connection to the named catalog.
  \item The four-script KLGHS pipeline: \texttt{scalarize.py}, \texttt{unify.py},
        \texttt{build\_chaos\_nulls.py}, \texttt{build\_pinch\_table.py}.
\end{itemize}
No proprietary databases, no specialist instruments, no institutional access
requirements. This is a deliberate constraint, not a limitation. A signal
that requires specialist infrastructure to observe is a signal that can be
attributed to the infrastructure. The KLGHS signal must survive the simplest
possible toolchain.

%% ============================================================
\section{Theoretical Basis: The Same-Object Principle}
\label{sec:same_object}

\subsection{The Kish Clutch and Register Families}

The KLGHS instrument transforms any physical measurement $x$ into a
dimensionless lattice scalar via the Kish Clutch operator:
\begin{equation}
  s(x) = \frac{\log(1 + x/x_0)}{\log(\kgeo)},
  \qquad \kgeo = \frac{16}{\pi} \approx 5.093
  \label{eq:kish_clutch}
\end{equation}
where $x_0$ is the domain-appropriate natural unit. The resulting scalar $s$ is
evaluated against the discrete harmonic register family:
\begin{equation}
  R_N = \frac{N}{\pi}, \quad N \in \{5, 6, 7, \ldots, 26\}
\end{equation}
Physical domains are not organized by scale. They are organized by physical
\emph{typology}. The Rosetta Atlas (Kish et al., 2026) establishes the
following empirically confirmed family taxonomy:

\begin{table}[h!]
\centering
\caption{Harmonic Family Taxonomy (derived from Volumes 5--10 confirmed results).}
\label{tab:rosetta_families}
\begin{tabular}{lll}
\toprule
\textbf{Family} & \textbf{Registers} & \textbf{Physical Type} \\
\midrule
Biological & $7/\pi$ & Organic geometry, amino acid backbone \\
Quantum / Chemical & $11/\pi$, $12/\pi$ & Bond lengths, atomic spectra, molecular vibration \\
Kinematic & $15/\pi$--$18/\pi$ & Velocity, tidal intervals, orbital periods, spin periods \\
Structural / Binding & $19/\pi$--$22/\pi$ & Binding energy, galaxy dispersion, nuclear decay \\
EM / Ceiling & $20/\pi$, $24/\pi$, $25/\pi$ & Colour, luminosity, CMB anisotropy \\
\bottomrule
\end{tabular}
\end{table}

\subsection{The Same-Object Principle}

The key insight motivating this programme is that most large public catalogs
contain measurements of \emph{multiple} physical attributes for the
\emph{same} objects. The confirmed results of Volume 11 already demonstrate
this: the same $\sim$1.8 million Gaia DR3 stars, measured four different ways,
yield four different registers:

\begin{center}
\begin{tabular}{llll}
\toprule
Lake & Attribute & Register & Z-score \\
\midrule
s1 & Parallax / distance & $24/\pi$ & (positional ceiling) \\
s2 & Radial velocity & $15/\pi$ & $z = +91.80$ \\
s3 & Absolute luminosity & $25/\pi$ & $z = +40.18$ \\
s4 & BP-RP colour & $20/\pi$ & $z = +135.20$ \\
\bottomrule
\end{tabular}
\end{center}

The register assignment follows the physical \emph{type} of the measurement, not
the objects being measured. The same stars express kinematic structure in
velocity, electromagnetic structure in colour and luminosity, and positional
structure in distance. This is not a coincidence of the dataset; it is a
statement about what kind of question is being asked.

\subsection{The Falsifiability Power of Same-Object Multi-Attribute Tests}

If $n$ attributes of the same object pool each cluster at their pre-registered
register \emph{family}, and each test is considered independent (conservative
lower bound), the probability of $n$ simultaneous confirmations under random
chance is:
\begin{equation}
  P(\text{all } n \text{ confirm by chance}) \approx \left(\frac{1}{5}\right)^n
\end{equation}
where $1/5$ is an upper bound on the probability of landing in the correct one
of five register families at random. For $n = 9$ attributes of the same 1.8M
stars: $P \lesssim 10^{-6}$. For $n = 9$ \emph{with the additional constraint
that each attribute's register is consistent with the Rosetta typology}, the
bound tightens considerably. This is not a fluke. It is a geometric fact about
the physical objects.

%% ============================================================
\section{Pre-Registered Predictions: Gaia DR3 Stellar Pool Extension}
\label{sec:gaia}

\subsection{Object Pool Definition}

The object pool for all s-series lake extensions is the Gaia DR3 Radial
Velocity Spectrometer (RVS) sample, filtered identically to the existing s1--s4
lakes:
\begin{itemize}
  \item Source: Gaia DR3, ESA Gaia Archive (\url{https://gea.esac.esa.int/archive/})
  \item Table: \texttt{gaiadr3.gaia\_source}
  \item Filter: \texttt{radial\_velocity IS NOT NULL AND parallax\_over\_error > 10
        AND ruwe < 1.4}
  \item Approximate $n$: 1.8 million stars
  \item Additional field extractions from same ADQL query
\end{itemize}

\subsection{New Lake Predictions}

% Corrected Column Widths for A4 Paper
\begin{longtable}{p{1.0cm} p{5.0cm} p{2.2cm} p{3.5cm} p{2.5cm} p{1.5cm}}
\caption{Gaia DR3 multi-attribute lake pre-registered predictions. All use
log-standard scalarization: $s = \log(1+x)/\log(16/\pi)$ unless noted.}
\label{tab:gaia_predictions}\\
\toprule
\textbf{Lake} & \textbf{Field / Description} & \textbf{Catalog Field} &
\textbf{Predicted Family} & \textbf{Predicted Register} & \textbf{Coverage}\\
\midrule
\endfirsthead
\multicolumn{6}{c}{\textit{(Table \ref{tab:gaia_predictions} continued)}}\\
\toprule
\textbf{Lake} & \textbf{Field / Description} & \textbf{Catalog Field} &
\textbf{Predicted Family} & \textbf{Predicted Register} & \textbf{Coverage}\\
\midrule
\endhead
\midrule
\multicolumn{6}{r}{\textit{Continued on next page}}\\
\endfoot
\bottomrule
\endlastfoot

s5 &
Proper motion magnitude $\sqrt{\mu_\alpha^2 + \mu_\delta^2}$ (mas/yr) &
\texttt{pmra}, \texttt{pmdec} &
Kinematic &
$15/\pi$ or $16/\pi$ &
$\sim$100\% \\[0.3em]

s6 &
Effective temperature $T_{\rm eff}$ (K) from GSP-Phot &
\texttt{teff\_gspphot} &
EM / Thermal &
$20/\pi$ or $25/\pi$ &
$\sim$87\% \\[0.3em]

s7 &
Surface gravity $\log g$ (cgs) from GSP-Phot &
\texttt{logg\_gspphot} &
Structural / Binding &
$19/\pi$--$22/\pi$ &
$\sim$87\% \\[0.3em]

s8 &
Metallicity $[\mathrm{M/H}]$ from GSP-Phot &
\texttt{mh\_gspphot} &
Quantum / Chemical &
$11/\pi$ or $12/\pi$ &
$\sim$75\% \\[0.3em]

s9 &
3D space velocity magnitude $\sqrt{v_r^2 + v_t^2}$ (km/s) &
Derived from \texttt{radial\_velocity}, \texttt{pmra}, \texttt{pmdec},
\texttt{parallax} &
Kinematic &
$15/\pi$ or $16/\pi$ &
$\sim$100\% \\[0.3em]

s10 &
Stellar radius $R_\star/R_\odot$ (FLAME pipeline) &
\texttt{radius\_flame} &
EM or Structural &
Open prediction &
$\sim$65\% \\[0.3em]

s11 &
Stellar mass $M_\star/M_\odot$ (FLAME pipeline) &
\texttt{mass\_flame} &
Structural / Binding &
$21/\pi$ or $22/\pi$ &
$\sim$65\% \\
\end{longtable}

\subsection{Rationale for Each Prediction}

\paragraph{s5 --- Proper motion (kinematic):}
Proper motion is the angular velocity of a star across the sky. It is a
kinematic quantity, measuring transverse motion relative to the line of sight,
distinct from radial velocity (s2) which measures motion along the line of
sight. The Rosetta Atlas places all velocity-type measurements in the kinematic
family (registers $15/\pi$--$18/\pi$). \textbf{Prediction: s5 will lock at the
same kinematic register as s2 ($15/\pi$ or $16/\pi$), demonstrating that two
independent measurements of stellar motion for the same stars confirm the same
register family.}

\paragraph{s6 --- Effective temperature (EM/thermal):}
$T_{\rm eff}$ quantifies the photospheric radiation field of the star ---
equivalent to the colour temperature of the emitted spectrum. It is an
electromagnetic quantity measuring how much energy per unit area the star
radiates. Colour (s4, confirmed at $20/\pi$) and $T_{\rm eff}$ are both
electromagnetic observables of the same physical process. \textbf{Prediction:
s6 will lock at the EM family ($20/\pi$ or $25/\pi$), demonstrating that two
independent EM measurements of the same stars confirm the same register family.}

\paragraph{s7 --- Surface gravity (structural):}
$\log g$ measures how strongly the stellar surface is gravitationally bound
--- the ratio of mass to radius-squared (in natural units). It is a structural
binding quantity. The Rosetta Atlas places binding resistance measurements in
the Structural/Binding family. \textbf{Prediction: s7 will lock in the
structural/binding register range ($19/\pi$--$22/\pi$). This is a new
register territory for stellar data and constitutes a novel prediction.}

\paragraph{s8 --- Metallicity (quantum/chemical):}
$[\mathrm{M/H}]$ is a dimensionless logarithmic ratio of metal abundance relative
to hydrogen, compared to solar composition. It encodes the chemical composition
of the stellar atmosphere at the atomic level --- a quantum/chemical quantity.
The Rosetta Atlas places chemistry, bond lengths, and atomic spectra at registers
$11/\pi$--$12/\pi$. \textbf{Prediction: s8 will lock at the quantum/chemical
register family ($11/\pi$ or $12/\pi$), demonstrating the framework detects
chemical composition structure at stellar scale.}

\paragraph{s9 --- 3D space velocity (kinematic):}
The full three-dimensional space velocity combines the line-of-sight component
(radial velocity, s2) with the plane-of-sky component (transverse velocity,
derived from proper motion and distance). It is the most complete kinematic
descriptor of stellar motion. \textbf{Prediction: s9 will lock at the kinematic
register ($15/\pi$ or $16/\pi$). Together with s2 and s5, this constitutes a
triple confirmation of kinematic register for the same 1.8M stars using three
geometrically independent velocity measurements.}

%% ============================================================
\section{Pre-Registered Predictions: Other Object Pools}
\label{sec:other_pools}

\subsection{ATNF Pulsar Catalog --- Multi-Attribute Extension}

\textbf{Object pool:} Same 2,527 pulsars as lake k2. Source: ATNF Pulsar
Catalogue (\url{https://www.atnf.csiro.au/research/pulsar/psrcat/}), accessed
via \texttt{psrcat} utility.

\begin{table}[h!]
\centering
\caption{ATNF pulsar multi-attribute predictions.}
\begin{tabular}{lp{4.5cm}p{4cm}l}
\toprule
\textbf{Lake} & \textbf{Field} & \textbf{Predicted Family} & \textbf{Register} \\
\midrule
k2a & Spin-down rate $\dot{P}$ (s/s) & Kinematic secondary & $15/\pi$--$18/\pi$ \\
k2b & Dispersion measure DM (pc/cm$^3$) & Structural / Positional & Open prediction \\
k2c & Characteristic age $P/(2\dot{P})$ (yr) & Kinematic / Temporal & $15/\pi$--$18/\pi$ \\
k2d & Surface B-field $B \propto \sqrt{P\dot{P}}$ (T) & Structural / Binding & $19/\pi$--$22/\pi$ \\
\bottomrule
\end{tabular}
\end{table}

\textbf{Key rationale:} The dispersion measure (k2b) is particularly diagnostic.
DM is an integral of free electron density along the line of sight --- a
completely different physical quantity from spin period. If k2b clusters at a
different register than k2, this confirms the framework assigns by physical
type, not by the objects observed. If it clusters at the same register, that is
equally informative and must be documented verbatim.

\subsection{NASA Exoplanet Archive --- Multi-Attribute Extension}

\textbf{Object pool:} Same 13,521 confirmed exoplanets as lakes p1, p2, p3.
Source: NASA Exoplanet Archive (\url{https://exoplanetarchive.ipac.caltech.edu/}).

\begin{table}[h!]
\centering
\caption{NASA exoplanet multi-attribute predictions.}
\begin{tabular}{lp{4.5cm}p{4cm}l}
\toprule
\textbf{Lake} & \textbf{Field} & \textbf{Predicted Family} & \textbf{Register} \\
\midrule
p\_ecc & Orbital eccentricity $e$ (dimensionless) & Structural shape & Null or $19/\pi$ \\
p\_incl & Orbital inclination $i$ (degrees) & Positional / geometric & Null predicted \\
p\_transit & Transit duration $T_{14}$ (hours) & Kinematic / temporal & $16/\pi$--$18/\pi$ \\
\bottomrule
\end{tabular}
\end{table}

\textbf{Note on transit duration:} Transit duration shares the same physical
units as tidal interval gaps (hours), which confirmed at $17/\pi$--$18/\pi$
across Atlantic, Gulf, Pacific, and Indian Ocean basins. If transit duration
also clusters in the kinematic family, this would extend the kinematic principle
from ocean tides to planetary orbital geometry.

\subsection{CERN Di-Muon --- Multi-Attribute Extension}

\textbf{Object pool:} Same 100,000 collision events as lake h1. Source: CMS
Open Data Portal, dataset \texttt{Run2012B\_DoubleMuParked}.

\begin{table}[h!]
\centering
\caption{CERN di-muon multi-attribute predictions.}
\begin{tabular}{lp{4.5cm}p{4cm}l}
\toprule
\textbf{Lake} & \textbf{Field} & \textbf{Predicted Family} & \textbf{Register} \\
\midrule
h1b & Leading muon $p_T$ (GeV/c) & Kinematic & $15/\pi$ or $16/\pi$ \\
h1c & Sub-leading muon $p_T$ (GeV/c) & Kinematic & $15/\pi$ or $16/\pi$ \\
h1d & Di-muon pseudorapidity $\eta$ & Positional / geometric & Null predicted \\
h1e & Relative azimuthal angle $\Delta\phi$ & Geometric & Open prediction \\
\bottomrule
\end{tabular}
\end{table}

\textbf{Significance of h1b and h1c:} Transverse momentum $p_T$ is a kinematic
quantity at sub-nuclear scale. If h1b/h1c confirm at the kinematic register
($15/\pi$--$16/\pi$), the kinematic principle extends from stellar velocities
(km/s, parsec scale) to particle momenta (GeV/c, femtometre scale) --- a span
of approximately 25 orders of magnitude for the kinematic register alone.

\subsection{SDSS DR16 Galaxies --- Multi-Attribute Extension}

\textbf{Object pool:} Same 1.843 million galaxies as lake g1. Source: SDSS
DR16 NASA-Sloan Atlas (\url{https://www.sdss.org/}).

\begin{table}[h!]
\centering
\caption{SDSS galaxy multi-attribute predictions.}
\begin{tabular}{lp{4.5cm}p{4cm}l}
\toprule
\textbf{Lake} & \textbf{Field} & \textbf{Predicted Family} & \textbf{Register} \\
\midrule
g1a & Effective radius $R_e$ (kpc) & EM / Ceiling & $24/\pi$ or $25/\pi$ \\
g1b & S\'{e}rsic index $n$ & Structural / morphological & $19/\pi$--$22/\pi$ \\
g1c & Stellar mass $M_\star$ ($M_\odot$) & Structural / Binding & $21/\pi$ \\
\bottomrule
\end{tabular}
\end{table}

\textbf{Significance of g1c:} Galaxy stellar mass and velocity dispersion are
connected empirically by the Faber-Jackson relation. If g1c (stellar mass)
confirms at the same register as g1 (velocity dispersion, $z = +62.89$ at
$21/\pi$), the KLGHS framework provides a geometric account of why the
Faber-Jackson relation exists: both quantities measure the same physical
property (resistance to structural dispersal) and therefore occupy the same
harmonic register.

%% ============================================================
\section{Dimensionality Falsification Protocol}
\label{sec:dim_falsification}

\subsection{The Scalarization Framework and Its Physical Meaning}

The KLGHS pipeline employs three distinct scalarization pathways, each encoding
a different physical relationship between a measurement and geometric tension:

% Corrected Column Widths for A4 Paper
\begin{table}[h!]
\centering
\caption{KLGHS scalarization pathways.}
\begin{tabular}{p{2.8cm} p{5.5cm} p{6cm}}
\toprule
\textbf{Transform} & \textbf{Formula} & \textbf{Physical Application} \\
\midrule
Log-standard & $s = \log(1+x)\,/\,\log(\kgeo)$ &
  Quantities directly proportional to geometric accumulation: distance, velocity,
  mass, duration, binding energy \\[0.5em]
Log-inverse & $s = \log(1+1/x)\,/\,\log(\kgeo)$ &
  Quantities inversely proportional to geometric tension: wavelength, period
  of high-energy transitions \\[0.5em]
Log-volumetric & $s = \log(1+x^{1/3})\,/\,\log(\kgeo)$ or
  $s = \log(1+x^3)\,/\,\log(\kgeo)$ &
  3D volumetric properties: bulk modulus, compressibility, material stiffness \\
\bottomrule
\end{tabular}
\end{table}

The choice of scalarization is not arbitrary. It is determined by the physical
dimensionality of the measurement: whether the attribute is intrinsically
one-dimensional (a speed, a duration, an energy), two-dimensional (an area,
a cross-section), three-dimensional (a volume, a bulk elastic property), or
an inverse quantity. Applying the wrong scalarization is equivalent to asking
the wrong physical question.

\subsection{The Empirical Precedent: The Materials Domain Collapse}
\label{sec:materials_collapse}

\textbf{This is not a hypothetical test. We have observed this effect.}

During the construction of the Materials domain lake (bulk modulus and
compressibility data, $n = 33{,}973$ records, confirmed at $19/\pi$,
$z = +69.77$), an initial processing error applied the 1D log-standard
scalarization to what is intrinsically a 3D volumetric property. The result was
immediate and unambiguous: the signal collapsed into the noise floor. No register
exceeded the $z > 3.0$ spectral identification threshold. The domain appeared
as a pure null.

When the correct volumetric scalarization was applied, the signal instantly
locked at $19/\pi$ with $z = +69.77$. This accidental discovery provides direct empirical evidence that:
\begin{enumerate}
  \item The KLGHS pipeline does not impose harmonic structure on all continuous
        data. It only detects structure that corresponds to the physical
        dimensionality supplied.
  \item The scalarization formula is a \emph{physical choice}, not a
        mathematical convenience. The wrong formula destroys the signal.
  \item The lattice geometry ``asks a specific question'' about the data. It
        only receives an answer when the question matches the physical nature
        of the attribute.
\end{enumerate}

In Volume 11, we formally document the Materials domain collapse under wrong
scalarization as a sovereign result, running it through the full MD5-verified
pipeline engine for inclusion in the official dataset record.

\subsection{Why This Test Is Scientifically Critical}

The strongest technically-competent version of the skeptical objection to the
KishLattice framework is the following:

\begin{quote}
\textit{``The log-modulo transform imposes harmonic structure on any sufficiently
large continuous dataset regardless of physical content. The observed z-scores
are artifacts of the mathematical pipeline rather than properties of the
physical data.''}
\end{quote}

This objection cannot be answered by showing higher z-scores. It must be
answered by demonstrating that the pipeline \emph{fails} when the physical
premise is wrong. A physical spectroscope --- such as a diffraction grating or a mass
spectrometer --- only reveals structure that is physically present in the input.
If you illuminate a diffraction grating with incoherent light, you see no
interference pattern. If you place the wrong sample in a mass spectrometer, you
read the wrong mass. The instrument does not manufacture signal from noise; it
reveals signal that was already there.

The Dimensionality Falsification Protocol establishes that KLGHS behaves as a
physical spectroscope:
\begin{itemize}
  \item Correct physical question $\rightarrow$ signal locked at predicted register.
  \item Wrong physical question $\rightarrow$ signal collapses to noise floor.
\end{itemize}

\textbf{This is the most direct possible answer to the transform-artifact
objection.} If the log-modulo mathematics were the source of the signal, it
would persist regardless of which scalarization formula is applied. If the
signal collapses under wrong dimensionality, the mathematics is not the source.
The physics is.

\subsection{Orthogonal Null Targets and Predictions}

We define four \emph{orthogonal null} tests. Each takes a confirmed strong
domain and intentionally applies the wrong scalarization formula.
\vspace{0.5em}

\paragraph{Orthogonal Null 1: Gaia Radial Velocity (s2\_wrong)}

\textbf{Confirmed domain:} s2, Gaia DR3 radial velocity, confirmed at $15/\pi$
with $z_{\rm real} = +91.80$.

\textbf{Correct scalarization:} Log-standard (1D kinematic):
\[
  s_{\rm correct} = \frac{\log(1 + v_r)}{\log(16/\pi)}
\]

\textbf{Intentional error --- volumetric scalarization applied to a 1D kinematic
quantity:}
\[
  s_{\rm wrong} = \frac{\log(1 + v_r^3)}{\log(16/\pi)}
\]

\textbf{Prediction:} $z_{\rm wrong} < 3.0$ at all registers. Complete phase
destruction.

\textbf{Falsification condition:} If $z_{\rm wrong} \geq 5.0$ at any register,
the transform-artifact hypothesis gains support and must be investigated before
submission.
\vspace{0.5em}

\paragraph{Orthogonal Null 2: ATNF Pulsar Spin Period (k2\_wrong)}

\textbf{Confirmed domain:} k2, ATNF pulsar spin period, confirmed STRONG in
kinematic family.

\textbf{Correct scalarization:} Log-standard (1D temporal/kinematic):
\[
  s_{\rm correct} = \frac{\log(1 + P)}{\log(16/\pi)}
\]

\textbf{Intentional error --- log-inverse (energy-type) scalarization applied
to a period:}
\[
  s_{\rm wrong} = \frac{\log(1 + 1/P)}{\log(16/\pi)}
\]

\textbf{Prediction:} $z_{\rm wrong} < 3.0$ at all registers. The spin period is
not an inverse-type quantity; the log-inverse formula has no physical
justification for this domain.
\vspace{0.5em}

\paragraph{Orthogonal Null 3: Exoplanet Orbital Period (p1\_wrong)}

\textbf{Confirmed domain:} p1, NASA exoplanet orbital periods, confirmed STRONG
at $22/\pi$, $z_{\rm real} \approx +78$.

\textbf{Correct scalarization:} Log-standard (1D temporal/kinematic):
\[
  s_{\rm correct} = \frac{\log(1 + P_{\rm orb})}{\log(16/\pi)}
\]

\textbf{Intentional error --- volumetric scalarization applied to a temporal
quantity:}
\[
  s_{\rm wrong} = \frac{\log(1 + P_{\rm orb}^3)}{\log(16/\pi)}
\]

\textbf{Prediction:} $z_{\rm wrong} < 3.0$ at all registers. An orbital period
in days is not a volumetric quantity; the cube transform has no dimensional
justification.

\textbf{Additional diagnostic:} The companion lake p2 (orbital radius, same
planets) is already an anti-signal. The contrast between p1\_correct (STRONG),
p1\_wrong (predicted NULL), and p2 (anti-signal) constitutes a complete
picture from a single planetary dataset.
\vspace{0.5em}

\paragraph{Orthogonal Null 4: Materials Domain Formal Documentation}

\textbf{Historical precedent:} During lake construction, bulk modulus data
was accidentally processed with 1D log-standard scalarization. Signal collapsed
to noise floor. Applying correct volumetric scalarization restored signal at
$19/\pi$, $z = +69.77$.

\textbf{Volume 11 action:} Run the incorrect 1D scalarization on the Materials
domain through the Vol 11 MD5-verified engine. Document the resulting
$z_{\rm wrong}$ vector formally in \texttt{z\_scores\_master.json}.

\textbf{Prediction:} All registers return $z < 3.0$ under the wrong
scalarization. This converts a historical observation into a sovereign pipeline
result with full chain of custody.

\subsection{Summary Table: Orthogonal Null Pre-Registrations}

% Corrected Column Widths for A4 Paper
\begin{table}[h!]
\centering
\caption{Orthogonal null pre-registered predictions.}
\begin{tabular}{l p{4.0cm} l p{4.5cm} l}
\toprule
\textbf{Lake} & \textbf{Source} & \textbf{Real $z$} &
\textbf{Wrong Formula} & \textbf{Predicted $z_{\rm wrong}$} \\
\midrule
s2\_wrong & Gaia radial velocity & $+91.80$ & $s = \log(1+v^3)/\log(\kgeo)$ & $< 3.0$ \\
k2\_wrong & ATNF spin period & STRONG & $s = \log(1+1/P)/\log(\kgeo)$ & $< 3.0$ \\
p1\_wrong & Exoplanet orbital period & $\sim+78$ & $s = \log(1+P^3)/\log(\kgeo)$ & $< 3.0$ \\
mat\_wrong & Materials bulk modulus & $+69.77$ & $s = \log(1+x)/\log(\kgeo)$ (1D) & $< 3.0$ \\
\bottomrule
\end{tabular}
\end{table}

%% ============================================================
\section{Complete Predictions Register}
\label{sec:predictions_register}

% Corrected Column Widths for A4 Paper
\begin{longtable}{p{1.5cm} p{5.0cm} p{3.0cm} p{3.5cm} p{1.5cm}}
\caption{Full pre-registration predictions table. All lakes to be built
\emph{after} Zenodo DOI assignment.}
\label{tab:full_predictions}\\
\toprule
\textbf{Lake ID} & \textbf{Physical Attribute} & \textbf{Source Catalog} &
\textbf{Predicted Family} & \textbf{Status} \\
\midrule
\endfirsthead
\multicolumn{5}{c}{\textit{(Table \ref{tab:full_predictions} continued)}}\\
\toprule
\textbf{Lake ID} & \textbf{Physical Attribute} & \textbf{Source Catalog} &
\textbf{Predicted Family} & \textbf{Status} \\
\midrule
\endhead
\midrule \multicolumn{5}{r}{\textit{Continued on next page}}\\
\endfoot
\bottomrule
\endlastfoot

\multicolumn{5}{l}{\textit{Gaia Stellar Pool Extensions}}\\[0.2em]
s5 & Proper motion magnitude (mas/yr) & Gaia DR3 & Kinematic ($15/\pi$--$16/\pi$) & \pending \\
s6 & $T_{\rm eff}$ (K) & Gaia DR3 GSP-Phot & EM / Thermal ($20/\pi$, $25/\pi$) & \pending \\
s7 & $\log g$ (cgs) & Gaia DR3 GSP-Phot & Structural ($19/\pi$--$22/\pi$) & \pending \\
s8 & $[\mathrm{M/H}]$ metallicity & Gaia DR3 GSP-Phot & Quantum/Chemical ($11/\pi$, $12/\pi$) & \pending \\
s9 & 3D space velocity (km/s) & Gaia DR3 derived & Kinematic ($15/\pi$--$16/\pi$) & \pending \\
s10 & Stellar radius ($R_\odot$) & Gaia DR3 FLAME & EM or Structural (open) & \pending \\
s11 & Stellar mass ($M_\odot$) & Gaia DR3 FLAME & Structural ($21/\pi$, $22/\pi$) & \pending \\[0.3em]

\multicolumn{5}{l}{\textit{ATNF Pulsar Extensions}}\\[0.2em]
k2a & Spin-down rate $\dot{P}$ (s/s) & ATNF & Kinematic secondary & \pending \\
k2b & Dispersion measure DM & ATNF & Structural / positional (open) & \pending \\
k2c & Characteristic age (yr) & ATNF derived & Kinematic / temporal & \pending \\
k2d & Surface B-field (T) & ATNF derived & Structural ($19/\pi$--$22/\pi$) & \pending \\[0.3em]

\multicolumn{5}{l}{\textit{Exoplanet Extensions}}\\[0.2em]
p\_ecc & Orbital eccentricity $e$ & NASA Exoplanet Archive & Null predicted & \pending \\
p\_incl & Inclination angle (deg) & NASA Exoplanet Archive & Null predicted & \pending \\
p\_transit & Transit duration (hr) & NASA Exoplanet Archive & Kinematic ($16/\pi$--$18/\pi$) & \pending \\[0.3em]

\multicolumn{5}{l}{\textit{CERN Di-Muon Extensions}}\\[0.2em]
h1b & Leading muon $p_T$ (GeV/c) & CMS Open Data & Kinematic ($15/\pi$, $16/\pi$) & \pending \\
h1c & Sub-leading muon $p_T$ (GeV/c) & CMS Open Data & Kinematic ($15/\pi$, $16/\pi$) & \pending \\
h1d & Di-muon pseudorapidity $\eta$ & CMS Open Data & Null predicted & \pending \\
h1e & Relative azimuthal angle $\Delta\phi$ & CMS Open Data & Open prediction & \pending \\[0.3em]

\multicolumn{5}{l}{\textit{SDSS Galaxy Extensions}}\\[0.2em]
g1a & Effective radius $R_e$ (kpc) & SDSS NSA & EM / Ceiling ($24/\pi$, $25/\pi$) & \pending \\
g1b & S\'{e}rsic index $n$ & SDSS NSA & Structural ($19/\pi$--$22/\pi$) & \pending \\
g1c & Stellar mass $M_\star$ ($M_\odot$) & SDSS NSA & Structural ($21/\pi$) & \pending \\[0.3em]

\multicolumn{5}{l}{\textit{Dimensionality Falsification (Orthogonal Nulls)}}\\[0.2em]
s2\_wrong & Gaia velocity, wrong (cubic) scalar & Gaia DR3 & $z < 3.0$ all registers & \pending \\
k2\_wrong & Pulsar period, wrong (inverse) scalar & ATNF & $z < 3.0$ all registers & \pending \\
p1\_wrong & Orbital period, wrong (cubic) scalar & NASA EPA & $z < 3.0$ all registers & \pending \\
mat\_wrong & Materials, wrong (1D) scalar & AFLOW / CSD & $z < 3.0$ all registers & \pending \\

\end{longtable}

%% ============================================================
\section{Falsification Conditions}
\label{sec:falsification}

The following conditions, if observed, would constitute a meaningful challenge
to the KLGHS framework and must be investigated before any submission:

\begin{enumerate}
  \item \textbf{Orthogonal null survives:} Any orthogonal null lake returns
  $z \geq 5.0$ at any register. This would indicate the signal persists
  regardless of dimensional correctness and supports the transform-artifact
  hypothesis.
  \item \textbf{Same-object register inconsistency:} Two kinematic attributes
  (e.g.\ s2 and s5) of the same 1.8M stars lock at different register
  \emph{families} (e.g.\ one at kinematic, one at structural). The register
  family taxonomy predicts they should agree; if they do not, the taxonomy
  requires revision.
  \item \textbf{Chaos null breach:} Any confirmed domain's z-score, on
  re-run through the Vol 11 engine, departs from its Vol 10 published value
  by more than the documented Monte Carlo variance band ($\pm 20\%$).
  \item \textbf{Systematic register drift:} If all new attribute lakes cluster
  at a single unexpected register regardless of physical type (e.g.\ all return
  $22/\pi$), this suggests a dataset-level artifact rather than a physical
  signal.
\end{enumerate}

A pre-registration that only lists expected confirmations is not a pre-registration
--- it is an announcement. The conditions above define the empirical boundaries
of the framework and are documented here in good faith.

%% ============================================================
\section{One-Lane Discipline: Why Sovereign Lakes Are Non-Negotiable}
\label{sec:one_lane}

One reviewer may ask: why not combine all attributes into a joint analysis?
This section documents the scientific rationale for maintaining strict
one-lake-one-question discipline.

\begin{enumerate}
  \item \textbf{Structural immunity to the cherry-picking objection.}
  If multiple attributes are combined into one lake, a referee can always ask:
  ``Did you try all combinations and report only the best?'' A Zenodo timestamp
  on each individual lake prediction, predating its construction, makes this
  objection structurally impossible. There is nothing to select after the fact
  when every selection is documented before the fact.

  \item \textbf{Minimum replication burden.}
  Any single sovereign lake can be independently replicated by any researcher
  with four hours, a laptop, and a public catalog. A merged multi-attribute lake
  requires understanding the merging logic, transform choices, and weighting
  scheme. Simplicity is the replication standard.

  \item \textbf{Storage cost is the price of honesty.}
  Separate lakes cost disk space and pipeline runtime. These costs are
  deliberately accepted. The benefit is that no one can question the process.

  \item \textbf{Null results are findings.}
  When p\_incl (orbital inclination) returns a null, that is a result. It
  documents that inclination angle is not a quantity that couples to the
  harmonic register structure. In a merged lake, a null attribute suppresses
  a confirming attribute and the null is invisible. In sovereign lakes, every
  outcome is on the record.

  \item \textbf{The pattern of confirmations and nulls is itself the finding.}
  Kinematic quantities confirm. Positional quantities produce anti-signal or
  null. EM quantities confirm at the EM ceiling. This taxonomy --- which domains
  light up and which do not --- is the empirical foundation of the Rosetta Atlas.
  It is only visible when each test is isolated.
\end{enumerate}

\textbf{We are being judged.} Every methodological decision in this programme
is made with the awareness that the results will face the highest level of
scrutiny. The one-lane discipline is not a constraint. It is the primary source
of the framework's credibility.

%% ============================================================
\section*{Acknowledgements}

The authors thank the ESA Gaia mission, the ATNF pulsar team, NASA for the
Exoplanet Archive, CERN for the CMS Open Data initiative, and the Sloan
Digital Sky Survey collaboration for maintaining the public catalogs that
underpin this programme. All data sources are publicly accessible without
institutional access requirements.

%% ============================================================
\section*{Data Availability}

All pipeline code, lake build scripts, and the MD5-verified unified master
dataset are available at:\\
\url{https://github.com/TimothyKish/Holographic-Resonance-The-Geometry-of-a-Quantized-Universe}

The pipeline requires Python 3.x with numpy, scipy, and pandas. No proprietary
libraries or databases are required.

%% ============================================================
\begin{thebibliography}{99}

\bibitem{rosetta2026}
Kish, T.J., Kish, P.A., Kish, L.A., \& Kish, M.A. (2026).
\textit{The Rosetta Atlas: The Master Equation of the Kish Lattice Theory.}
Zenodo. \url{https://doi.org/10.5281/zenodo.19939949}

\bibitem{vol10}
Kish, T.J., Kish, P.A., Kish, L.A., \& Kish, M.A. (2026).
\textit{KishLattice Geometric Harmonic Spectroscopy: The Deep Survey.}
Zenodo. \url{https://doi.org/10.5281/zenodo.20279208}

\bibitem{vol9}
Kish, T.J. et al. (2026).
\textit{KishLattice Geometric Harmonic Spectroscopy: The First Survey of a New Field.}
Zenodo. \url{https://doi.org/10.5281/zenodo.19935604}

\bibitem{vol9prereg}
Kish, T.J. \& Kish, M.A. (2026).
\textit{Volume 9 --- Predictions Pre-Registration.}
Zenodo. \url{https://doi.org/10.5281/zenodo.19702023}

\bibitem{gaia_dr3}
Gaia Collaboration (2023).
\textit{Gaia Data Release 3: Summary of the content and survey properties.}
\textit{Astronomy \& Astrophysics}, 674, A1.

\bibitem{atnf}
Manchester, R.N. et al. (2005).
\textit{The Australia Telescope National Facility Pulsar Catalogue.}
\textit{The Astronomical Journal}, 129, 1993.

\end{thebibliography}

\end{document}